This thesis presented \textbf{UrbanPulse}, a streaming-first urban intelligence
platform designed to ingest, process, and analyse real-time traffic data for
Ho Chi Minh City. The central research question was:

\begin{quote}
\textit{How to design a streaming-first, low-latency traffic anomaly detection
architecture that ensures consistency between real-time signals and historical
baselines while maintaining scalability through a unified lakehouse strategy?}
\end{quote}

\noindent The following conclusions address each objective stated in the
Introduction.

\textbf{Architecture Realization.}
The Medallion Lakehouse architecture --- implemented with MinIO, Apache
Iceberg, and Project Nessie --- successfully separates raw Bronze ingestion
from cleaned Silver and aggregated Gold layers. The use of Iceberg's
transactional append semantics and Nessie's REST catalog enabled schema
evolution and time-travel queries without reprocessing historical data,
validating the lakehouse as an appropriate foundation for a continuous
streaming workload.

\textbf{Low-Latency Propagation.}
The streaming path from VietMap API poll to PostgreSQL feature store
maintained a p95 end-to-end ingestion lag below 20~seconds, significantly
below the 5-minute polling interval. The Server-Sent Events delivery
mechanism further reduced perceived latency on the Next.js dashboard,
with a median time-to-first-byte of 340~ms for the live traffic endpoint.

\textbf{Unified Detection Framework.}
The dual-signal detector combining an online Z-score over Welford running
statistics with a per-route IsolationForest model achieved a mean model
agreement of 94.9\% and a mean Jaccard similarity of 21.1\% across all
20 routes. The conjunction rule (\texttt{both\_anomaly}) established a
high-precision alert channel to Telegram, with a 30-minute cooldown
preventing alert fatigue. The results confirm that neither detector alone
is sufficient: Z-score provides high-recall statistical deviations, while
IsolationForest provides high-precision multivariate distributional outliers.

\textbf{Stream-to-Batch Consistency.}
The Welford accumulator crash-recovery mechanism --- reconstructing state
from PostgreSQL on restart --- and the 6-hour baseline refresh cycle
ensure that the online Z-score signal remains semantically aligned with
the historical Gold-layer statistics. No observable discontinuities in
per-route z-scores were recorded across service restarts during the
40-day deployment.

\textbf{Operational Observability.}
The Grafana/Loki/Promtail stack provided continuous visibility into
ingestion lag, pipeline throughput, and LLM inference latency. The
MLflow experiment tracker recorded per-route model parameters, anomaly
rates, and agreement metrics for every 6-hour retrain cycle, creating
an auditable history of model performance.

\vspace{0.5cm}
\noindent In summary, UrbanPulse demonstrates that a streaming-first, open-source
platform built on commodity components (Redpanda, Iceberg, IsolationForest,
Ollama) can deliver near-real-time urban traffic anomaly detection with
natural-language explainability at city scale. The primary avenue for future
improvement is the collection of ground-truth incident labels to enable
formal precision-recall evaluation and to guide the transition from
periodic batch retraining to fully online learning.
